\subsection{Idempotents and Semisimple Rings}

Let A be a ring. Recall that an element e of A is called idempotent (I, §1, No.~4, p.~7) if we have \(e^{2} = e\) . It is then also an idempotent element of the opposite ring \(A^{o}\) of A.

PROPOSITION 13. --- a) A left ideal \(\mathfrak{a}\) of A admits a supplement in \(\mathrm{A}_s\) if and only if there exists an idempotent \(e\) in A such that \(\mathfrak{a} = \mathrm{Ae}\) . The ideal \(\mathfrak{a}\) then consists of the elements \(x\) of A such that \(x = xe\) .

\begin{enumerate}
\def\labelenumi{\alph{enumi})}
\setcounter{enumi}{1}
\item
  Let \(e\) and \(f\) be idempotents in A. We have \(Ae \subset Af\) if and only if we have \(ef = e\) .
\item
  Let M be an A-module. Then M is projective and monogenous if and only if there exists an idempotent e in A such that M is isomorphic to Ae.
\end{enumerate}

The endomorphisms of the A-module \(A_{s}\) are the right multiplications by the elements of A. The projectors in the A-module \(A_{s}\) are therefore the mappings \(x \mapsto xe\) where e is an idempotent in A. Moreover, the submodules of \(A_{s}\) are the left ideals, and such a submodule admits a supplement if and only if it is the image of a projector (II, §1, No.~9, p.~211, Proposition 14). Assertion a) follows.

The relation \(\mathrm{Ae} \subset \mathrm{Af}\) is equivalent to \(e \in \mathrm{Af}\) . By a), it is therefore equivalent to \(e = ef\) ; assertion b) follows.

If the A-module M is monogenous, there exists a surjective A-linear mapping \(u: A_{s} \to M\) . If, moreover, M is projective, then there exists a submodule a of \(A_{s}\) supplementary to the kernel of u. Then u induces an isomorphism from a to M. Conversely, if M is isomorphic to a direct factor of \(A_{s}\) , then it is monogenous and projective. Assertion c) therefore follows from a).

Remarks. --- 1) Let \(\mathfrak{a}\) be a left ideal of A. By the proof above and the corollary of Proposition 12 of II, §1, No.~8, p.~209, the mapping \(e \mapsto \mathrm{A}(1 - e)\) defines a bijection from the set of idempotents \(e\) in A such that \(\mathfrak{a} = \mathrm{A}e\) to the set of left ideals \(\mathfrak{b}\) of A such that \(\mathrm{A}_s = \mathfrak{a} \oplus \mathfrak{b}\) .

\begin{enumerate}
\def\labelenumi{\arabic{enumi})}
\setcounter{enumi}{1}
\tightlist
\item
  Let \(e\) and \(f\) be idempotents in A. By Proposition 13, b), we have \(Ae = Af\) if and only if \(ef = e\) and \(fe = f\) . Consequently, if the ring A is commutative, then the relation \(Ae = Af\) is equivalent to \(e = f\) . This does not hold in general, as shown by the example \(A = M_2(Z)\) , \(e = \begin{pmatrix} 1 & 0 \\ 0 & 0 \end{pmatrix}\) , and \(f = \begin{pmatrix} 1 & 0 \\ 1 & 0 \end{pmatrix}\) .
\end{enumerate}

We say that idempotents \(e\) and \(e'\) in the ring A are orthogonal if \(ee' = e'e = 0\) . Let \((e_i)_{i \in I}\) be a finite family of pairwise orthogonal idempotents in A. Since we have

\[
\left(\sum_ {i} e _ {i}\right) ^ {2} = \sum_ {i} e _ {i} ^ {2} + \sum_ {i \neq j} e _ {i} e _ {j} = \sum_ {i} e _ {i},
\]

the element \(\sum_{i\in I}e_i\) of A is idempotent.

A partition of an idempotent e in A is a finite family \((e_{i})_{i\in\mathrm{I}}\) of pairwise orthogonal idempotents in A such that \(e=\sum_{i\in I}e_{i}\) . We say that an idempotent e in A is decomposable if there exists a partition of e consisting of pairwise orthogonal idempotents distinct from 0 and e; in the opposite case, we say that it is indecomposable. Observe that 0 is a decomposable idempotent.

PROPOSITION 14. --- Let e be an idempotent in A.

\begin{enumerate}
\def\labelenumi{\alph{enumi})}
\item
  If \((e_i)_{i\in I}\) is a partition of \(e\) , then the A-module Ae is the direct sum of the family \((\mathrm{A}e_i)_{i\in I}\) .
\item
  Let \((\mathfrak{a}_{i})_{i\in\mathrm{I}}\) be a finite family of left ideals of A with direct sum Ae. For \(i\in I\) , denote the component of e in \(a_{i}\) by \(e_{i}\) . Then \((e_{i})_{i\in\mathrm{I}}\) is a partition of e, and we have \(a_{i}=Ae_{i}\) for every \(i\in I\) .
\item
  The A-module Ae is indecomposable if and only if the idempotent e is indecomposable.
\end{enumerate}

Let \((e_i)_{i\in \mathrm{I}}\) be a partition of \(e\) . For every \(i\in \mathrm{I}\) , we have

\[
e _ {i} e = \sum_ {j} e _ {i} e _ {j} = e _ {i} ^ {2} + \sum_ {j \neq i} e _ {i} e _ {j} = e _ {i},
\]

and therefore \(Ae_{i} \subset Ae\) . For every \(i \in I\) , we define an A-linear projector \(p_{i}\) in Ae by setting \(p_{i}(x) = xe_{i}\) . We have \(p_{i}p_{j} = 0\) if \(i \neq j\) , and for every x in Ae,

\[
x = x e = \sum_ {i} x e _ {i} = \sum_ {i} p _ {i} (x).
\]

Consequently (II, §1, No.~8, p.~20, Proposition 12), Ae is the direct sum of the images of the \(p_i\) . Now, we have \(ee_i = e_i e = e_i\) , so the image of \(p_i\) is \(\mathrm{Ae}_i\) . This proves a).

Take the notation and assumptions of b). Let \(i \in I\) . Since \(e_{i}\) belongs to Ae, we have \(e_{i} = e_{i}e = \sum_{j} e_{i}e_{j}\) . Since Ae is the direct sum of the \(a_{j}\) and \(e_{i}e_{j}\) belongs to \(a_{j}\) for every j, we have \(e_{i} = e_{i}e_{i}\) and \(e_{i}e_{j} = 0\) for \(i \neq j\) . In other words, \((e_{i})_{i \in I}\) is a partition of the idempotent e. By a), Ae is the direct sum of the \(Ae_{i}\) . By assumption, we have \(Ae_{i} \subset a_{i}\) , and Ae is the direct sum of the \(a_{i}\) . We therefore have \(Ae_{i} = a_{i}\) for every \(i \in I\) . We have proved b).

Finally, c) follows immediately from a) and b).

Remark 3. --- Applying the previous results to the opposite ring of A, we see, in particular, that a monogenous right A-module M is projective if and only if there exists an idempotent e in A such that M is isomorphic to eA. Moreover, eA is an indecomposable right module if and only if e is indecomposable.

Now suppose that the ring A is semisimple, and denote the set of classes of simple left A-modules by S. The A-module \(A_{s}\) is semisimple, and every submodule of \(A_{s}\) is a direct factor. Let a be a left ideal of A. By the above, there exists an idempotent e in A such that a = Ae, and the A-module a is simple if and only if it is indecomposable, that is, if and only if e is indecomposable.

Let \((\mathfrak{m}_i)_{i\in \mathrm{I}}\) be a family of minimal left ideals of A such that we have \(\mathrm{A}_s = \oplus_{i\in \mathrm{I}}\mathfrak{m}_i\) . By Proposition 14, there exists a partition \((\varepsilon_i)_{i\in \mathrm{I}}\) of 1 consisting of indecomposable idempotents such that \(\mathfrak{m}_i = \mathrm{A}\varepsilon_i\) for every \(i\in \mathrm{I}\) .

For every \(\lambda \in \mathcal{S}\) , let \(S_{\lambda}\) be an A-module of class \(\lambda\) , and let \(\mathfrak{a}_{\lambda}\) be the isotypical component of type \(\lambda\) of the A-module \(A_s\) . Since A is the direct sum of the family \((\mathfrak{a}_{\lambda})_{\lambda \in \mathcal{S}}\) , there exists a partition \((e_{\lambda})_{\lambda \in \mathcal{S}}\) of 1 such that \(\mathfrak{a}_{\lambda} = Ae_{\lambda}\) for every \(\lambda \in \mathcal{S}\) . For \(\lambda \in \mathcal{S}\) , denote by I( \(\lambda\) ) the set of indices \(i \in I\) such that the simple A-module \(\mathfrak{m}_i\) is of type \(\lambda\) . By Proposition 4, b) of VIII, p.~65, we have

\[
\mathfrak {a} _ {\lambda} = \bigoplus_ {i \in I (\lambda)} \mathfrak {m} _ {i}.\tag{3}
\]

The idempotent \(e_{\lambda}\) is the component of 1 in \(\mathfrak{a}_{\lambda}\) , so \(e_{\lambda} = \sum_{i\in \mathrm{I}(\lambda)}\varepsilon_i\) .

PROPOSITION 15. --- Suppose that the ring A is semisimple.

\begin{enumerate}
\def\labelenumi{\alph{enumi})}
\item
  For every \(\lambda \in \mathcal{S}\) , \(e_{\lambda}\) is the unique element of the center \(Z\) of \(A\) satisfying the relations \((e_{\lambda})_{S_{\lambda}} = 1_{S_{\lambda}}\) and \((e_{\lambda})_{S_{\mu}} = 0\) for \(\mu \neq \lambda\) .
\item
  The indecomposable idempotents in the ring Z are the \(e_{\lambda}\) , and the minimal ideals of Z are the \(Ze_{\lambda}\) for \(\lambda \in S\) .
\item
  Let M be an A-module and \((\mathrm{M}_{\lambda})_{\lambda\in\mathcal{S}}\) the family of its isotypical components. The family of projectors associated with the decomposition of M into the direct sum of the \(M_{\lambda}\) (VIII, p.~65) is \(((e_{\lambda})_{\mathrm{M}})_{\lambda\in\mathcal{S}}\) , and we have \(M_{\lambda}=a_{\lambda}M\) for every \(\lambda\in S\) .
\end{enumerate}

Let \(\lambda\) and \(\mu\) be distinct in \(\mathcal{S}\) . We have \(e_{\lambda} \in \mathfrak{a}_{\lambda}\) , and \(\mathfrak{a}_{\lambda}\) is contained in the annihilator \(\mathfrak{b}_{\mu}\) of the A-module \(\mathrm{S}_{\mu}\) (VIII, p.~142, Proposition 9). We therefore have \((e_{\lambda})_{\mathrm{S}_{\mu}} = 0\) . The relation \((e_{\lambda})_{\mathrm{S}_{\lambda}} = 1_{\mathrm{S}_{\lambda}}\) follows because we have \(1 = \sum_{\nu \in \mathcal{S}} e_{\nu}\) . Assertion a) is then a consequence of Proposition 8 of VIII, p.~141.

Let \(\lambda\) be in \(\mathcal{S}\) . The two-sided ideal \(\mathfrak{a}_{\lambda}\) of \(A\) consists of the elements \(x\) such that \(x = xe_{\lambda}\) ; we therefore have \(Z \cap \mathfrak{a}_{\lambda} = Ze_{\lambda}\) , and therefore b) by Proposition 10, b).

Let us prove c). Let \(x\) be an element of M. We have \(e_{\lambda} \in \mathfrak{a}_{\lambda}\) for every \(\lambda \in \mathcal{S}\) . Since the mapping \(a \mapsto ax\) from \(A_s\) to M is A-linear, we have \(\mathfrak{a}_{\lambda}x \subset M_{\lambda}\) (VIII, p.~66, Proposition 5) and, in particular, \(e_{\lambda}x \in M_{\lambda}\) . We have \(1 = \sum_{\lambda \in \mathcal{S}} e_{\lambda}\) , hence \(x = \sum_{\lambda \in \mathcal{S}} e_{\lambda}x\) . Consequently, \(e_{\lambda}x\) is the component of \(x\) in \(M_{\lambda}\) .

Remark 4. --- Suppose that the ring A is semisimple. Let \((e_{i})_{i\in\mathrm{I}}\) be a partition of 1 consisting of nonzero idempotents in the center Z of A. If \(\operatorname{Card}(\mathrm{I})=\operatorname{Card}(\mathcal{S})\) , then the \(e_{i}\) are the indecomposable idempotents in Z.

